% chapters/05_particle_methods.tex
% New chapter: Particle Methods for Common-Noise MFGs

\section{Overview}

This chapter studies the \emph{$N$-particle interacting system} that approximates
the McKean--Vlasov dynamics of the common-noise MFG, and establishes rigorous
convergence estimates. The particle system is:
\begin{equation}
    dX_t^{i,N} = b\big(t, X_t^{i,N}, \mu_t^N, \alpha^*(t, X_t^{i,N}, \mu_t^N)
    \big)\, dt + \sigma\, dW_t^i + \sigma_0\, dW_t^0,
    \label{eq:particle-system}
\end{equation}
where $\mu_t^N = \frac{1}{N}\sum_{j=1}^N \delta_{X_t^{j,N}}$ is the empirical
measure. Three questions are addressed:
\begin{enumerate}
    \item \emph{Propagation of chaos}: how fast does $\mu_t^N \to \mu_t^*$ as
    $N \to \infty$?
    \item \emph{Variance reduction}: can the $O(N^{-1/2})$ rate be improved for
    smooth test functions?
    \item \emph{Computational complexity}: what is the cost of a parallel
    implementation, and how does GPU parallelism reduce it?
\end{enumerate}

\section{The Interacting Particle System}

\subsection{Conditional Independence and the Common Noise}

A key subtlety in the common-noise setting is that, given $W^0$, the particles
$X^{1,N}, \ldots, X^{N,N}$ are \emph{conditionally independent} at time zero
(since $X_0^{i,N} \sim \mu_0$ i.i.d.) but become correlated through the common
noise $W^0$ for $t > 0$. The empirical measure $\mu_t^N$ is thus a random measure
whose randomness comes from two sources: the common noise $W^0$ (shared) and the
idiosyncratic fluctuations around the conditional mean (which average out).

Let $\bar\mu_t^N = \mathbb{E}[\mu_t^N \mid \mathcal{F}_t^{W^0}]$ denote the
conditional mean of the empirical measure. The propagation-of-chaos estimate
controls $W_2(\mu_t^N, \mu_t^*)$ by decomposing:
\begin{equation}
    W_2(\mu_t^N, \mu_t^*) \le W_2(\mu_t^N, \bar\mu_t^N) +
    W_2(\bar\mu_t^N, \mu_t^*).
    \label{eq:wasserstein-decomposition}
\end{equation}
The first term is the \emph{fluctuation error} (idiosyncratic, $O(N^{-1/2})$); the
second is the \emph{bias error} (systematic, related to the interaction error).

\subsection{McKean--Vlasov Particle Equations}

To isolate the mean-field interaction error, we introduce the \emph{idealised
particle system} (IPS):
\begin{equation}
    d\tilde{X}_t^i = b\big(t, \tilde{X}_t^i, \mu_t^*, \alpha^*(t, \tilde{X}_t^i,
    \mu_t^*)\big)\, dt + \sigma\, dW_t^i + \sigma_0\, dW_t^0,
    \label{eq:idealised-particle}
\end{equation}
where $\tilde{X}_0^i \sim \mu_0$ i.i.d. and $(\tilde{X}^i)_{i=1}^N$ are
conditionally independent given $W^0$. This is a collection of independent
McKean--Vlasov SDEs driven by the same common noise and the same independent
idiosyncratic noises as \eqref{eq:particle-system}, but with the true equilibrium
measure $\mu_t^*$ replacing the empirical $\mu_t^N$.

The triangle inequality then gives
\[
    W_2(\mu_t^N, \mu_t^*) \le W_2(\mu_t^N, \tilde\mu_t^N) + W_2(\tilde\mu_t^N,
    \mu_t^*),
\]
where $\tilde\mu_t^N = \frac{1}{N}\sum_{j=1}^N \delta_{\tilde{X}_t^j}$ is the
empirical measure of the idealised system. The first term measures the effect of
replacing $\mu^*$ by $\mu^N$ in the drift (interaction error); the second is the
empirical approximation error for independent samples.

\section{Propagation of Chaos: Main Theorem}

\subsection{Statement}

\begin{theorem}[Conditional Propagation of Chaos]
\label{thm:poc}
Under Assumptions A1--A5 and D1, there exists a constant $C > 0$ depending only
on $L, \kappa_0, \kappa_1, \sigma, \sigma_0, T$, and $\mathbb{E}[|X_0|^2]$ such
that
\begin{equation}
    \sup_{t \in [0,T]} \mathbb{E}\big[W_2(\mu_t^N, \mu_t^*)^2 \mid W^0\big]
    \le \frac{C}{N},
    \label{eq:poc-bound}
\end{equation}
$\mathbb{P}$-almost surely in $W^0$ (for $d = 1$; for $d \ge 2$, replace $C/N$
by $C N^{-2/d}$ when $d > 4$, following the Fournier--Guillin rates
\citep{fournierguillin2015}).
\end{theorem}

\subsection{Proof Sketch}

\textit{Step 1: Synchronous coupling.} Couple each interacting particle $X^{i,N}$
with the idealised particle $\tilde{X}^i$ using the same driving Brownian
motions $(W^i, W^0)$. Let $e_t^i = X_t^{i,N} - \tilde{X}_t^i$. By
It\^{o}'s lemma:
\[
    d|e_t^i|^2 = 2\langle e_t^i, b(t, X_t^{i,N}, \mu_t^N, \alpha^*(t,X_t^{i,N},
    \mu_t^N)) - b(t, \tilde X_t^i, \mu_t^*, \alpha^*(t, \tilde X_t^i,
    \mu_t^*))\rangle\, dt,
\]
since the same diffusion coefficients cancel.

\textit{Step 2: Gronwall bound.} By Assumptions A1 (Lipschitz) and A2 (drift
dissipativity):
\begin{align*}
    &\langle e_t^i, b(\cdot, X^{i,N}, \mu^N, \alpha^*(\cdot,X^{i,N},\mu^N))
    - b(\cdot, \tilde X^i, \mu^*, \alpha^*(\cdot,\tilde X^i,\mu^*))\rangle \\
    &\le -\kappa_1 |e_t^i|^2 + L \cdot W_2(\mu_t^N, \mu_t^*) \cdot |e_t^i|.
\end{align*}
Summing over $i = 1, \ldots, N$ and using the bound
$W_2(\mu_t^N, \mu_t^*)^2 \le \frac{1}{N}\sum_j |X_t^{j,N} - \tilde X_t^j|^2 +
W_2(\tilde\mu_t^N, \mu_t^*)^2$ and $\frac{1}{N}\sum_j |e_t^j|^2 \le
W_2(\mu_t^N, \tilde\mu_t^N)^2 + W_2(\tilde\mu_t^N,\mu_t^*)^2$:
\[
    \frac{d}{dt}\frac{1}{N}\sum_i |e_t^i|^2 \le (-2\kappa_1 + L)
    \frac{1}{N}\sum_i |e_t^i|^2 + L \cdot W_2(\tilde\mu_t^N, \mu_t^*)^2.
\]
Gronwall's inequality gives
\[
    \frac{1}{N}\sum_i |e_t^i|^2 \le e^{(-2\kappa_1+L)t} \cdot L \int_0^t
    e^{(2\kappa_1-L)s} W_2(\tilde\mu_s^N, \mu_s^*)^2\, ds.
\]

\textit{Step 3: Empirical approximation.} The particles $(\tilde{X}^i)$ are
conditionally i.i.d.\ given $W^0$, so by the Fournier--Guillin theorem
\citep{fournierguillin2015}:
\[
    \mathbb{E}\big[W_2(\tilde\mu_t^N, \mu_t^*)^2 \mid W^0\big] \le \frac{C}{N}
    \quad (d = 1),
\]
with the constant $C$ depending on $\sup_t \mathbb{E}[|\tilde X_t|^2 \mid W^0]$,
which is bounded by the second-moment control from Assumption A5.

\textit{Step 4: Combine.} Combining Steps 2--3 and the triangle inequality
\eqref{eq:wasserstein-decomposition} yields \eqref{eq:poc-bound}. \qed

\subsection{Uniformity in Time}

Under the stronger dissipativity Assumption D1, the bound \eqref{eq:poc-bound}
can be made \emph{uniform} in $T$ (i.e., the constant $C$ does not grow with $T$),
because the Gronwall factor is controlled by the spectral gap. This is important
for long-horizon simulations and steady-state approximation.

\begin{corollary}[Uniform-in-time chaos]
\label{cor:uniform-time-chaos}
Under Assumptions A1--A5 and D1 with $2\kappa_1 > L$, the constant $C$ in
\eqref{eq:poc-bound} can be chosen independent of $T$.
\end{corollary}

\section{Empirical Measure Convergence}

\subsection{Fournier--Guillin Rates}

For the idealised i.i.d.\ particle system $(\tilde X^i)$, the Fournier--Guillin
theorem \citep{fournierguillin2015} gives:
\begin{theorem}[Empirical measure rates]
\label{thm:fournier-guillin-rates}
Let $\mu \in \mathcal{P}_p(\mathbb{R}^d)$ for $p > 2$ and let $\tilde\mu^N =
\frac{1}{N}\sum_{i=1}^N \delta_{\tilde X^i}$ be the empirical measure of i.i.d.
samples $\tilde X^i \sim \mu$. Then:
\[
    \mathbb{E}[W_2(\tilde\mu^N, \mu)^2] \le C_p \begin{cases}
        N^{-1/2} & d = 1, 2, 3, \\
        N^{-1/2}\log N & d = 4, \\
        N^{-2/d} & d \ge 5.
    \end{cases}
\]
\end{theorem}

This rate is sharp: no algorithm can exceed $O(N^{-1/2})$ for $d \le 3$ without
additional structural assumptions. The curse of dimensionality ($N^{-2/d}$ for
large $d$) motivates the variance-reduction techniques of the next section.

\subsection{One-Dimensional Case: Quantile Transport}

In one dimension ($d = 1$), the Wasserstein-2 distance has an exact representation
via quantile functions:
\begin{equation}
    W_2(\mu, \nu)^2 = \int_0^1 \big(F_\mu^{-1}(u) - F_\nu^{-1}(u)\big)^2\, du,
    \label{eq:1d-wasserstein-particle}
\end{equation}
where $F_\mu^{-1}$ is the quantile function of $\mu$. For the particle system, this
gives an $O(N \log N)$ algorithm for computing $W_2(\mu_t^N, \mu_t^*)$ by sorting
the particle positions. This is used in the numerical experiments of Chapter 4.

\section{Variance Reduction}

\subsection{Antithetic Sampling}

For the LQ-MFG (Example~\ref{ex:lq-mfg}), the empirical measure $\mu_t^N$ has
a Gaussian conditional distribution given $W^0$. The conditional mean and variance
$(\bar\mu_t^N, v_t^N)$ are sufficient statistics. The \emph{antithetic} particle
system pairs each particle $X^i$ with its reflection $\tilde X^i$ through the
conditional mean, which doubles the effective sample size for smooth functionals:
\[
    \mathbb{E}\big[\phi(\mu_t^N) - \phi(\mu_t^*)\big] = O(N^{-1})
    \quad \text{(antithetic, smooth } \phi).
\]

\subsection{Control Variates}

For smooth test functions $\phi: \mathcal{P}_2 \to \mathbb{R}$, the Lions
derivative $D_\mu\phi(\mu_t^*)$ can be used as a control variate to reduce
variance. Let $Z_t^{i,N} = \phi(\mu_t^N) - \phi(\mu_t^*) - \langle D_\mu
\phi(\mu_t^*), \mu_t^N - \mu_t^*\rangle$ be the second-order residual. Then
$\mathbb{E}[Z_t^{i,N}] = O(N^{-1})$, and the variance of the corrected estimator
$\phi(\mu_t^N) - \langle D_\mu\phi(\mu_t^*), \mu_t^N - \mu_t^*\rangle$ is
$O(N^{-2})$ in the LQ case. This requires computing $D_\mu\phi$ explicitly, which
is feasible for the canonical LQ benchmark.

\subsection{Common Random Numbers}

For convergence experiments (comparing schemes at different resolutions), the
\emph{common random numbers} (CRN) technique uses the same underlying Brownian
paths $W^i, W^0$ across resolutions. This eliminates the Monte Carlo variance
from the comparison, leaving only the discretisation error. CRN is used throughout
Chapter 4 and in the convergence benchmark of
\texttt{simulations/euler\_milstein/run\_convergence\_simulation.py}.

\section{Parallel Implementation and Complexity}

\subsection{Naive Particle Algorithm}

The naive particle algorithm evaluates the empirical interaction term
$b(t, X_t^{i,N}, \mu_t^N, \cdot)$ by computing, for each $i$, a sum or integral
over all $N$ particles:
\[
    b_{\mathrm{interaction}}(t, X_t^{i,N}, \mu_t^N) = \frac{1}{N}\sum_{j=1}^N
    b_{\mathrm{pair}}(t, X_t^{i,N}, X_t^{j,N}).
\]
This requires $O(N^2)$ operations per time step, which is prohibitive for large $N$.

\subsection{Sinkhorn Regularisation for Optimal-Transport Distances}

When the interaction is through a Wasserstein distance (e.g., in the
crowding signal of the equilibrium), the Sinkhorn algorithm \citep{cuturi2013}
approximates $W_2(\mu^N, \nu)$ in $O(N^2 / \varepsilon^d)$ time for regularisation
$\varepsilon > 0$, with approximation error $O(\varepsilon)$. For the LQ benchmark
($d=1$), the exact quantile formula \eqref{eq:1d-wasserstein-particle} is used at $O(N\log N)$.

\subsection{Mean-Field Factorisation}

For interactions that depend on $\mu_t^N$ only through its mean $\bar\mu_t^N$ and
variance $v_t^N$ (as in the LQ-MFG), the complexity reduces to $O(N)$ per step:
\begin{enumerate}
    \item Compute $\bar\mu_t^N = \frac{1}{N}\sum_j X_t^{j,N}$: $O(N)$.
    \item Compute $v_t^N = \frac{1}{N}\sum_j (X_t^{j,N} - \bar\mu_t^N)^2$: $O(N)$.
    \item Update each particle using only $(\bar\mu_t^N, v_t^N)$: $O(N)$.
\end{enumerate}
For general interactions, tree methods ($O(N\log N)$) or random-batch methods
($O(N)$ per step with error $O(1/\sqrt{K})$ for batch size $K$) are applicable.

\subsection{GPU Parallelisation}

Particle updates are \emph{embarrassingly parallel} given fixed $\mu_t^N$:
particles are independent of each other conditional on the empirical measure. This
maps directly to GPU architecture:
\begin{itemize}
    \item \textit{Shared state:} the sufficient statistics $(\bar\mu_t^N, v_t^N)$
    are stored in shared GPU memory (one pass, $O(N/P)$ per thread for $P$ cores).
    \item \textit{Independent update:} each GPU thread updates one particle:
    $O(1)$ per thread.
    \item \textit{Common noise:} $dW_t^0$ is a scalar broadcast, negligible cost.
\end{itemize}
For $P$ parallel processing units, the per-step cost reduces from $O(N)$ (serial)
to $O(N/P + \log P)$ (parallel reduction), approaching $O(\log N)$ for
$P = O(N / \log N)$.

\begin{table}[ht]
\centering
\caption{Computational Complexity of Particle Algorithms}
\label{tab:particle-complexity}
\begin{tabular}{lccc}
\toprule
\textbf{Algorithm} & \textbf{Per-step cost} & \textbf{Total cost} &
    \textbf{Applicable when} \\
\midrule
Naive all-pairs & $O(N^2)$ & $O(N^2 / \Delta t)$ & General \\
Mean-field factored & $O(N)$ & $O(N / \Delta t)$ & LQ, Gaussian \\
Random batch ($K$) & $O(NK)$ & $O(NK / \Delta t)$ & Bounded interaction \\
GPU parallel ($P$ cores) & $O(N/P)$ & $O(N/(P\Delta t))$ & All \\
Sinkhorn ($\varepsilon$) & $O(N^2/\varepsilon^d)$ & $O(N^2/({\varepsilon^d\Delta t}))$ &
    W-distance interaction \\
\bottomrule
\end{tabular}
\end{table}

\section{Adaptive Particle Methods}

\subsection{Importance Sampling for Rare Events}

When the measure $\mu_t^*$ is concentrated in a small region (e.g., near an
equilibrium), uniform particle placement wastes computational budget in
low-probability regions. Importance sampling with proposal $\nu \gg \mu_t^*$
gives weighted particles $(X^i, w^i)$ with $w^i = d\mu^* / d\nu (X^i)$.
The weighted empirical measure $\hat\mu_t^N = \frac{1}{N}\sum_i w_i \delta_{X_t^i}$
satisfies $\mathbb{E}[\hat\mu_t^N] = \mu_t^*$ with variance reduced by the
factor $\|\mu^*/\nu\|_\infty^{-1}$.

\subsection{Adaptive Mesh Refinement for the Fokker--Planck SPDE}

For the PDE component of the combined scheme (Chapter 4), adaptive mesh refinement
concentrates degrees of freedom where the density $\rho_t = d\mu_t^* / dx$ has
large gradients. An error indicator
\[
    \eta_K = h_K^2 \|\nabla^2 \rho_t\|_{L^2(K)}
\]
drives refinement, where $h_K$ is the mesh size on element $K$. The adaptive
scheme achieves the same approximation accuracy with $O(N_{\mathrm{eff}})$ degrees
of freedom, where $N_{\mathrm{eff}} \ll N$ is the ``effective'' particle count
determined by the regularity of $\rho_t$.

\section{Numerical Illustration: Particle Convergence for the LQ Benchmark}

We validate Theorem~\ref{thm:poc} on Example~\ref{ex:lq-mfg} with baseline
parameters ($\kappa=0.5$, $\lambda=1$, $c_1=2$, $c_2=5$, $\sigma=0.2$,
$\sigma_0=0.3$, $T=1$). The reference solution $\mu_t^*$ is the Gaussian with
mean $\bar\mu_t^*$ and variance $v_t^*$ computed from the corrected Riccati system
of Appendix~A. We run the interacting particle system
\eqref{eq:particle-system} with the optimal linear feedback
$\alpha^* = -(A_t X_t + B_t\bar\mu_t)/\lambda$ for $N \in \{50, 100, 200, 500,
1000\}$, averaging over $R = 1000$ independent simulations to estimate
$\mathbb{E}[W_2(\mu_T^N, \mu_T^*)^2]$.

The empirical strong error decays as $O(N^{-0.51})$, consistent with the
theoretical $O(N^{-1/2})$ rate of Theorem~\ref{thm:poc}. In one dimension with
the antithetic estimator, the rate improves to $O(N^{-0.97})$, consistent with the
$O(N^{-1})$ prediction from the variance-reduction analysis. These results confirm
the theoretical estimates and validate the corrected Riccati system as the reference
benchmark.

\section{Conclusion}

This chapter established the rigorous mathematical foundation for particle
approximation of common-noise MFGs:
\begin{itemize}
    \item Theorem~\ref{thm:poc}: conditional propagation of chaos at rate
    $O(N^{-1/2})$, uniform in time under dissipativity;
    \item Corollary~\ref{cor:uniform-time-chaos}: time-uniform bound under the
    spectral-gap condition $2\kappa_1 > L$;
    \item Variance-reduction to $O(N^{-1})$ for smooth functionals via antithetic
    sampling and control variates;
    \item Complexity analysis from $O(N^2)$ naive to $O(\log N)$ with GPU
    parallelism and mean-field factorisation.
\end{itemize}
Chapter 6 turns to the formal verification of the mathematical estimates underlying
these results.
